\documentclass[11pt,a4paper]{article}
\usepackage[italian]{babel}
\usepackage[utf8]{inputenc}
\usepackage[T1]{fontenc}
\usepackage{amsmath,amssymb,amsthm}
\usepackage{geometry}
\geometry{margin=2.2cm}
\usepackage{enumitem}
\usepackage{tikz}
\usepackage{tabularx,array,booktabs}
\usepackage{xcolor}
\usepackage{hyperref}
\hypersetup{colorlinks=true, linkcolor=blue, urlcolor=blue}
\DeclareMathOperator{\arccot}{arccot}
\newcommand{\risp}{\vspace{2pt}\noindent\textbf{Soluzione.}\ }
\newcolumntype{Y}{>{\raggedright\arraybackslash}X}
\setlist[enumerate,1]{label=\alph*)}
\newcommand{\assi}{%
\draw[step=0.5,gray!20,very thin] (-4,-2) grid (4,3.5);
\draw[->] (-4,0) -- (4.2,0) node[right]{\small $x$};
\draw[->] (0,-2) -- (0,3.7) node[above]{\small $y$};
\draw[dashed,gray] (-2,-2) -- (3.5,3.5);}

\title{\textbf{Svolgimento} \\ \large Scheda di lavoro: le funzioni goniometriche inverse \\ Matematica -- Classe quarta}
\author{Liceo Scientifico ``Galileo Galilei'' -- San Don\`a di Piave}
\date{}

\begin{document}
\maketitle
\vspace{-1cm}
\noindent\textit{Documento riservato al docente: contiene lo svolgimento completo della scheda.}

\section*{Parte A -- Grafici}

\subsection*{A.1 -- A.2 Tabelle}
\begin{center}
\renewcommand{\arraystretch}{1.5}
\begin{tabular}{|c|*{7}{c|}}
\hline
$x$ & $-\frac{\pi}{2}$ & $-\frac{\pi}{3}$ & $-\frac{\pi}{6}$ & $0$ & $\frac{\pi}{6}$ & $\frac{\pi}{3}$ & $\frac{\pi}{2}$\\\hline
$\sin x$ & $-1$ & $-\frac{\sqrt3}{2}$ & $-\frac12$ & $0$ & $\frac12$ & $\frac{\sqrt3}{2}$ & $1$\\\hline
\end{tabular}
\\[6pt]
\begin{tabular}{|c|*{7}{c|}}
\hline
$x$ & $0$ & $\frac{\pi}{6}$ & $\frac{\pi}{3}$ & $\frac{\pi}{2}$ & $\frac{2\pi}{3}$ & $\frac{5\pi}{6}$ & $\pi$\\\hline
$\cos x$ & $1$ & $\frac{\sqrt3}{2}$ & $\frac12$ & $0$ & $-\frac12$ & $-\frac{\sqrt3}{2}$ & $-1$\\\hline
\end{tabular}
\end{center}
Scambiando le righe si ottengono, per esempio, $\arcsin\frac{\sqrt3}{2}=\frac{\pi}{3}$ e $\arccos\left(-\frac12\right)=\frac{2\pi}{3}$.

\noindent\textbf{Perch\'e non $\left[-\frac{\pi}{2},\frac{\pi}{2}\right]$ per il coseno?} Perch\'e l\`i il coseno non \`e iniettivo ($\cos\left(-\frac{\pi}{3}\right)=\cos\frac{\pi}{3}=\frac12$) e non assume valori negativi, quindi non \`e suriettivo su $[-1,1]$. In $[0,\pi]$ il coseno \`e strettamente decrescente e assume tutti i valori di $[-1,1]$.

\subsection*{A.3 Tabelle}
\begin{center}
\renewcommand{\arraystretch}{1.5}
\begin{tabular}{|c|*{5}{c|}}
\hline
$x$ & $-\frac{\pi}{3}$ & $-\frac{\pi}{4}$ & $0$ & $\frac{\pi}{4}$ & $\frac{\pi}{3}$\\\hline
$\tan x$ & $-\sqrt3$ & $-1$ & $0$ & $1$ & $\sqrt3$\\\hline
\end{tabular}
\\[6pt]
\begin{tabular}{|c|*{5}{c|}}
\hline
$x$ & $\frac{\pi}{6}$ & $\frac{\pi}{4}$ & $\frac{\pi}{2}$ & $\frac{3\pi}{4}$ & $\frac{5\pi}{6}$\\\hline
$\cot x$ & $\sqrt3$ & $1$ & $0$ & $-1$ & $-\sqrt3$\\\hline
\end{tabular}
\end{center}
Gli asintoti verticali $x=\pm\frac{\pi}{2}$ della tangente diventano gli asintoti orizzontali $y=\pm\frac{\pi}{2}$ dell'arcotangente; per l'arcocotangente gli asintoti $x=0$ e $x=\pi$ della cotangente diventano $y=0$ e $y=\pi$.

\subsection*{Grafici}
In blu la funzione goniometrica ristretta, in rosso la sua inversa, tratteggiata la bisettrice $y=x$.

\begin{center}
\begin{tikzpicture}[scale=0.62]
\assi
\draw[thick,blue,domain=-1.5708:1.5708,samples=60] plot (\x,{sin(\x r)});
\draw[thick,red,domain=-1:1,samples=80] plot (\x,{rad(asin(\x))});
\node[red] at (-2.3,2.9) {\small $y=\arcsin x$};
\node[blue] at (2.6,-1.3) {\small $y=\sin x$};
\end{tikzpicture}
\qquad
\begin{tikzpicture}[scale=0.62]
\assi
\draw[thick,blue,domain=0:3.1416,samples=60] plot (\x,{cos(\x r)});
\draw[thick,red,domain=-1:1,samples=80] plot (\x,{rad(acos(\x))});
\draw[dotted] (-4,1.5708) -- (4,1.5708) node[right]{\scriptsize $\frac{\pi}{2}$};
\node[red] at (-2.4,2.9) {\small $y=\arccos x$};
\node[blue] at (2.8,-1.4) {\small $y=\cos x$};
\end{tikzpicture}
\\[8pt]
\begin{tikzpicture}[scale=0.62]
\assi
\draw[thick,blue,domain=-1.1:1.1,samples=60] plot (\x,{tan(\x r)});
\draw[thick,red,domain=-4:4,samples=80] plot (\x,{rad(atan(\x))});
\draw[dashed,red] (-4,1.5708) -- (4,1.5708) node[right]{\scriptsize $\frac{\pi}{2}$};
\draw[dashed,red] (-4,-1.5708) -- (4,-1.5708) node[right]{\scriptsize $-\frac{\pi}{2}$};
\node[red] at (-2.6,2.6) {\small $y=\arctan x$};
\end{tikzpicture}
\qquad
\begin{tikzpicture}[scale=0.62]
\assi
\draw[thick,blue,domain=0.3:2.66,samples=60] plot (\x,{cos(\x r)/sin(\x r)});
\draw[thick,red,domain=-4:4,samples=80] plot (\x,{pi/2-rad(atan(\x))});
\draw[dashed,red] (-4,3.1416) -- (4,3.1416) node[right]{\scriptsize $\pi$};
\node[red] at (-2.3,1.2) {\small $y=\operatorname{arccot} x$};
\end{tikzpicture}
\end{center}

\subsection*{A.4 Tabella riassuntiva}
\begin{center}
\renewcommand{\arraystretch}{1.6}
\small
\begin{tabular}{|l|c|c|c|c|}
\hline
 & $y=\arcsin x$ & $y=\arccos x$ & $y=\arctan x$ & $y=\arccot x$\\\hline
Dominio & $[-1,1]$ & $[-1,1]$ & $\mathbb{R}$ & $\mathbb{R}$\\\hline
Codominio & $\left[-\frac{\pi}{2},\frac{\pi}{2}\right]$ & $[0,\pi]$ & $\left(-\frac{\pi}{2},\frac{\pi}{2}\right)$ & $(0,\pi)$\\\hline
Monotonia & crescente & decrescente & crescente & decrescente\\\hline
Simmetrie & dispari & centro $\left(0,\frac{\pi}{2}\right)$ & dispari & centro $\left(0,\frac{\pi}{2}\right)$\\\hline
Asintoti & nessuno & nessuno & $y=\pm\frac{\pi}{2}$ & $y=0$, $y=\pi$\\\hline
\end{tabular}
\end{center}
Relazioni utili: $\arccos(-x)=\pi-\arccos x$, \ $\arcsin x+\arccos x=\frac{\pi}{2}$, \ $\arccot x=\frac{\pi}{2}-\arctan x$, \ $\arccot(-x)=\pi-\arccot x$.

\section*{Parte B -- Valori e composizioni}
\subsection*{B.1}
\begin{enumerate}
\item $\arcsin\frac12=\frac{\pi}{6}$
\item $\arcsin\left(-\frac{\sqrt3}{2}\right)=-\frac{\pi}{3}$ \quad (disparit\`a)
\item $\arccos\left(-\frac{\sqrt2}{2}\right)=\pi-\frac{\pi}{4}=\frac{3\pi}{4}$
\item $\arccos 0=\frac{\pi}{2}$
\item $\arctan\left(-\sqrt3\right)=-\frac{\pi}{3}$
\item $\arctan 1=\frac{\pi}{4}$
\item $\arccot(-1)=\pi-\frac{\pi}{4}=\frac{3\pi}{4}$ \quad (non $-\frac{\pi}{4}$, che non sta in $(0,\pi)$)
\item $\arccot\sqrt3=\frac{\pi}{6}$
\end{enumerate}

\subsection*{B.2}
\begin{enumerate}
\item $\sin\left(\arcsin\frac13\right)=\frac13$, perch\'e $\sin(\arcsin x)=x$ per ogni $x\in[-1,1]$.
\item $\sin\frac{5\pi}{6}=\frac12$, quindi $\arcsin\left(\sin\frac{5\pi}{6}\right)=\arcsin\frac12=\frac{\pi}{6}\neq\frac{5\pi}{6}$, perch\'e $\frac{5\pi}{6}\notin\left[-\frac{\pi}{2},\frac{\pi}{2}\right]$.
\item $\cos\left(-\frac{\pi}{3}\right)=\frac12$, quindi il risultato \`e $\arccos\frac12=\frac{\pi}{3}$.
\item $\tan\frac{3\pi}{4}=-1$, quindi il risultato \`e $\arctan(-1)=-\frac{\pi}{4}$.
\item Posto $\alpha=\arcsin\frac35\in\left[-\frac{\pi}{2},\frac{\pi}{2}\right]$, si ha $\cos\alpha\ge0$, quindi $\cos\alpha=\sqrt{1-\frac{9}{25}}=\frac45$.
\item $\arccos\left(-\frac12\right)=\frac{2\pi}{3}$, quindi il risultato \`e $\tan\frac{2\pi}{3}=-\sqrt3$.
\item Non ha senso: $2\notin[-1,1]$, dominio dell'arcoseno.
\item Posto $\alpha=\arctan 2$, si ha $\alpha\in\left(0,\frac{\pi}{2}\right)$ perch\'e $2>0$. Da $\sin\alpha=\dfrac{\tan\alpha}{\sqrt{1+\tan^2\alpha}}$ (valida per $\cos\alpha>0$) segue $\sin\alpha=\dfrac{2}{\sqrt5}=\dfrac{2\sqrt5}{5}$.
\end{enumerate}

\section*{Parte C -- Domini}
\begin{enumerate}
\item $-1\le 2x-1\le 1 \iff 0\le x\le 1$. \quad $D=[0,1]$.
\item $-1\le x^2-1\le 1 \iff 0\le x^2\le 2 \iff -\sqrt2\le x\le\sqrt2$. \quad $D=\left[-\sqrt2,\sqrt2\right]$.
\item L'arcotangente \`e definita su tutto $\mathbb{R}$: basta $x\neq 2$. \quad $D=\mathbb{R}\setminus\{2\}$.
\item Serve $x\in[-1,1]$ e $\arcsin x\ge0$, cio\`e $0\le x\le 1$. \quad $D=[0,1]$.
\item Basta che esista $\sqrt x$: \quad $D=[0,+\infty)$.
\item Condizioni: $x\neq1$ e $-1\le\dfrac{x+1}{x-1}\le1$.
\[
\frac{x+1}{x-1}\le 1 \iff \frac{2}{x-1}\le 0 \iff x<1;
\qquad
\frac{x+1}{x-1}\ge -1 \iff \frac{2x}{x-1}\ge 0 \iff x\le 0 \ \vee\ x>1.
\]
Intersecando: $D=(-\infty,0]$.
\end{enumerate}

\section*{Parte D -- Grafici deducibili}
\begin{enumerate}
\item Traslazione di $\frac{\pi}{2}$ verso l'alto di $y=\arcsin x$: $D=[-1,1]$, codominio $[0,\pi]$, crescente, passa per $\left(0,\frac{\pi}{2}\right)$.\\
\emph{Osservazione:} coincide con il grafico di $y=\arccos x$ ribaltato rispetto all'asse $y$, cio\`e $\arcsin x+\frac{\pi}{2}=\arccos(-x)$. Infatti $\arccos(-x)=\pi-\arccos x=\pi-\left(\frac{\pi}{2}-\arcsin x\right)$.
\item Dilatazione verticale di fattore $2$: $D=\mathbb{R}$, codominio $(-\pi,\pi)$, asintoti orizzontali $y=\pm\pi$.
\item Traslazione di $1$ verso destra: $D=[0,2]$, codominio $[0,\pi]$, decrescente; passa per $(0,\pi)$, $\left(1,\frac{\pi}{2}\right)$, $(2,0)$.
\item Si ribalta rispetto all'asse $x$ la parte con $x<0$: $D=[-1,1]$, codominio $\left[0,\frac{\pi}{2}\right]$, funzione pari, minimo nell'origine con punto angoloso.
\end{enumerate}

\section*{Parte E -- Equazioni e problema}
\subsection*{E.1}
\begin{enumerate}
\item $\frac{\pi}{3}\in\left[-\frac{\pi}{2},\frac{\pi}{2}\right]$, quindi $x=\sin\frac{\pi}{3}=\frac{\sqrt3}{2}$.
\item $\frac{2\pi}{3}\in[0,\pi]$, quindi $2x=\cos\frac{2\pi}{3}=-\frac12$, da cui $x=-\frac14$.
\item $x-1=\tan\left(-\frac{\pi}{4}\right)=-1$, da cui $x=0$.
\item Impossibile: $-\frac{\pi}{6}\notin[0,\pi]$, codominio dell'arcocoseno.
\item Da $\arcsin x+\arccos x=\frac{\pi}{2}$ e $\arcsin x=\arccos x$ segue $\arcsin x=\frac{\pi}{4}$, cio\`e $x=\frac{\sqrt2}{2}$.
\end{enumerate}

\subsection*{E.2 La rampa}
\begin{enumerate}
\item La pendenza \`e il rapporto tra dislivello e percorso orizzontale, cio\`e la tangente dell'angolo di inclinazione $\alpha$. Quindi
\[
\alpha_{\max}=\arctan 0{,}08\approx 0{,}0798\ \text{rad}\approx 4{,}57^\circ.
\]
\item No: i $5$ m sono misurati lungo il piano inclinato, cio\`e sono l'ipotenusa. Il rapporto $\frac{0{,}40}{5}=0{,}08$ \`e il \emph{seno} dell'angolo:
\[
\alpha=\arcsin 0{,}08\approx 0{,}0801\ \text{rad}\approx 4{,}59^\circ>4{,}57^\circ.
\]
La pendenza effettiva \`e $\tan\alpha\approx 0{,}0803=8{,}03\%$ (proiezione orizzontale $5\cos\alpha\approx 4{,}98$ m), quindi la rampa \textbf{non \`e a norma}, anche se di poco.\\
\emph{Nota per il docente:} l'errore tipico \`e concludere ``$40/500=8\%$, a norma'', confondendo seno e tangente. Per angoli piccoli i due valori sono quasi uguali, e questo apre una discussione utile sull'approssimazione $\sin\alpha\approx\tan\alpha\approx\alpha$.
\end{enumerate}

\subsection*{E.3 Dimostrazione}
Sia $x\in[-1,1]$ e $\alpha=\arcsin x$, quindi $-\frac{\pi}{2}\le\alpha\le\frac{\pi}{2}$ e $\sin\alpha=x$. Allora $0\le\frac{\pi}{2}-\alpha\le\pi$ e $\cos\left(\frac{\pi}{2}-\alpha\right)=\sin\alpha=x$. L'angolo $\frac{\pi}{2}-\alpha$ sta nel codominio dell'arcocoseno e ha coseno $x$; poich\'e il coseno \`e iniettivo in $[0,\pi]$, \`e proprio $\arccos x$. Quindi $\arccos x=\frac{\pi}{2}-\arcsin x$, cio\`e la tesi.

\end{document}
